\documentclass[12pt]{article}
\usepackage{amsmath, amssymb, geometry}
\usepackage[T1]{fontenc}
\usepackage{geometry} % OK, already present in your document

\geometry{margin=1in}

% Global fixes for overfull boxes
\setlength{\emergencystretch}{1em}
\sloppy

\begin{document}
	
	\title{SVM Slides}
	\author{}
	\date{}
	\maketitle


\subsection*{Page 1 — Introduction}

\textbf{Content extract}
\begin{itemize}
	\item Introduction to Support Vector Machines (SVM).
	\item Part of the course \emph{Machine Learning I}.
\end{itemize}

\textbf{Exam-style questions}

\begin{enumerate}
	\item What is the primary task of Support Vector Machines?
	\begin{enumerate}
		\item Clustering
		\item Dimensionality reduction
		\item Classification
		\item Reinforcement learning
		\item Feature extraction
	\end{enumerate}
	\textbf{Answer: (c)}
	
	\item Are SVMs supervised or unsupervised learning methods?
	\begin{enumerate}
		\item Unsupervised
		\item Supervised
		\item Semi-supervised
		\item Reinforcement-based
		\item Self-supervised
	\end{enumerate}
	\textbf{Answer: (b)}
	
	\item What type of output does a basic SVM produce?
	\begin{enumerate}
		\item Continuous probability
		\item Cluster index
		\item Binary label
		\item Ranking score
		\item Regression value
	\end{enumerate}
	\textbf{Answer: (c)}
	
	\item True or False: SVMs were originally designed for regression.
	\begin{enumerate}
		\item True
		\item False
	\end{enumerate}
	\textbf{Answer: (b)}
	
	\item In which academic field do SVMs belong?
	\begin{enumerate}
		\item Computer Graphics
		\item Databases
		\item Machine Learning
		\item Operating Systems
		\item Software Engineering
	\end{enumerate}
	\textbf{Answer: (c)}
\end{enumerate}

\subsection*{Page 2 — Origin and Motivation}

\textbf{Content extract}
\begin{itemize}
	\item SVMs are derived from Statistical Learning Theory.
	\item Introduced by Vapnik and Chervonenkis in 1992.
\end{itemize}

\textbf{Exam-style questions}

\begin{enumerate}
	\item SVMs are grounded in which theory?
	\begin{enumerate}
		\item Bayesian Theory
		\item Statistical Learning Theory
		\item Information Theory
		\item Game Theory
		\item Graph Theory
	\end{enumerate}
	\textbf{Answer: (b)}
	
	\item Who introduced Support Vector Machines?
	\begin{enumerate}
		\item Hinton and LeCun
		\item Vapnik and Chervonenkis
		\item Bishop and Jordan
		\item Pearl and Judea
		\item Rumelhart and McClelland
	\end{enumerate}
	\textbf{Answer: (b)}
	
	\item When were SVMs introduced?
	\begin{enumerate}
		\item 1985
		\item 1990
		\item 1992
		\item 1998
		\item 2005
	\end{enumerate}
	\textbf{Answer: (c)}
	
	\item What is a major advantage of SVMs?
	\begin{enumerate}
		\item Low computational cost
		\item No need for labeled data
		\item Strong theoretical guarantees
		\item Automatic feature extraction
		\item No hyperparameters
	\end{enumerate}
	\textbf{Answer: (c)}
	
	\item True or False: SVMs always outperform neural networks.
	\begin{enumerate}
		\item True
		\item False
	\end{enumerate}
	\textbf{Answer: (b)}
\end{enumerate}

\subsection*{Page 3 — Training Data}

\textbf{Content extract}
\begin{itemize}
	\item Each training example is a pair $(\mathbf{x}_i, y_i)$.
	\item $\mathbf{x}_i \in \mathbb{R}^n$ is a feature vector.
	\item $y_i \in \{+1, -1\}$ is the class label.
\end{itemize}

\begin{enumerate}
	\item What does $\mathbf{x}_i$ represent?
	\begin{enumerate}
		\item A class label
		\item A feature vector
		\item A probability
		\item A margin
	\end{enumerate}
	\textbf{Answer: (b)}
	
	\item True or False: SVM labels are continuous values.
	
	\textbf{Answer: False}
	
	\item If $\mathbf{x}_i \in \mathbb{R}^5$, how many features does sample $i$ have?
	
	\textbf{Answer: 5}
\end{enumerate}


\subsection*{Page 4 — Hyperplane Definition}

\textbf{Content extract}
\begin{itemize}
	\item A hyperplane is defined as $\mathbf{w}^\top \mathbf{x} + b = 0$.
	\item $\mathbf{w}$ controls orientation, $b$ controls offset.
\end{itemize}

\begin{enumerate}
	\item What does $b$ control?
	\begin{enumerate}
		\item Margin width
		\item Orientation
		\item Offset of hyperplane
		\item Dimensionality
	\end{enumerate}
	\textbf{Answer: (c)}
	
	\item True or False: $\mathbf{w}$ is orthogonal to the hyperplane.
	
	\textbf{Answer: True}
	
	\item If $\mathbf{w} = (2,1)$ and $b=0$, write the hyperplane equation.
	
	\textbf{Answer: $2x_1 + x_2 = 0$}
\end{enumerate}


\subsection*{Pages 5--7 — Classification Rule}

\textbf{Content extract}
\begin{itemize}
	\item Decision function: $f(\mathbf{x}) = \mathbf{w}^\top \mathbf{x} + b$.
	\item Class is determined by the sign of $f(\mathbf{x})$.
\end{itemize}

\begin{enumerate}
	\item When is a sample classified as $+1$?
	\begin{enumerate}
		\item When $f(\mathbf{x}) < 0$
		\item When $f(\mathbf{x}) = 0$
		\item When $f(\mathbf{x}) > 0$
		\item Always
	\end{enumerate}
	\textbf{Answer: (c)}
	
	\item True or False: $f(\mathbf{x})=0$ lies on the decision boundary.
	
	\textbf{Answer: True}
	
	\item If $f(\mathbf{x})=-2.7$, what is the predicted class?
	
	\textbf{Answer: $-1$}
\end{enumerate}



\subsection*{Pages 8--11 — Multiple Hyperplanes}

\textbf{Content extract}
\begin{itemize}
	\item Many hyperplanes may separate the same training data.
	\item They differ in generalization performance.
\end{itemize}

\begin{enumerate}
	\item Why are not all separating hyperplanes equivalent?
	\begin{enumerate}
		\item Same margin
		\item Same error
		\item Different generalization
		\item Different labels
	\end{enumerate}
	\textbf{Answer: (c)}
	
	\item True or False: Any separating hyperplane is optimal.
	
	\textbf{Answer: False}
	
	\item How many hyperplanes can separate linearly separable data?
	
	\textbf{Answer: Infinitely many}
\end{enumerate}



\subsection*{Pages 12--17 — Margin Motivation}

\textbf{Content extract}
\begin{itemize}
	\item Hyperplanes too close to data are unstable.
	\item Noise may cause misclassification.
\end{itemize}

\begin{enumerate}
	\item What problem arises with small margins?
	\begin{enumerate}
		\item Underfitting
		\item Instability
		\item High bias
		\item No solution
	\end{enumerate}
	\textbf{Answer: (b)}
	
	\item True or False: Larger margins improve robustness.
	
	\textbf{Answer: True}
	
	\item If margin doubles, what happens to robustness?
	
	\textbf{Answer: It increases}
\end{enumerate}


\subsection*{Pages 18--20 — Margin Definition}

\textbf{Content extract}
\begin{itemize}
	\item Margin is the distance between class boundaries.
	\item SVM maximizes the margin.
\end{itemize}

\begin{enumerate}
	\item What does SVM optimize?
	\begin{enumerate}
		\item Training error
		\item Likelihood
		\item Margin width
		\item Distance to origin
	\end{enumerate}
	\textbf{Answer: (c)}
	
	\item True or False: All points affect the margin equally.
	
	\textbf{Answer: False}
	
	\item Which samples define the margin?
	
	\textbf{Answer: Support vectors}
\end{enumerate}


\subsection*{Pages 21--25 — Margin and Support Vectors}

\textbf{Content extract}
\begin{itemize}
	\item Margin is the distance between the two class-separating hyperplanes.
	\item The optimal hyperplane maximizes this margin.
	\item The simplest SVM is the Linear SVM (LSVM).
\end{itemize}

\begin{enumerate}
	\item What is the margin in SVM?
	\begin{enumerate}
		\item Distance between centroids
		\item Distance between class means
		\item Distance between separating hyperplanes
		\item Distance from origin
	\end{enumerate}
	\textbf{Answer: (c)}
	
	\item True or False: A larger margin improves generalization.
	
	\textbf{Answer: True}
	
	\item What type of SVM is obtained for perfectly separable data?
	
	\textbf{Answer: Linear (hard-margin) SVM}
\end{enumerate}


\subsection*{Pages 26--29 — Support Vectors}

\textbf{Content extract}
\begin{itemize}
	\item Support vectors are the samples that lie on the margin boundaries.
	\item Only support vectors determine the optimal hyperplane.
	\item Other samples can be ignored once the solution is found.
\end{itemize}

\begin{enumerate}
	\item Which samples define the decision boundary?
	\begin{enumerate}
		\item All samples
		\item Only misclassified samples
		\item Support vectors
		\item Random samples
	\end{enumerate}
	\textbf{Answer: (c)}
	
	\item True or False: Removing non-support vectors changes the classifier.
	
	\textbf{Answer: False}
	
	\item How many support vectors are required at minimum in $\mathbb{R}^2$?
	
	\textbf{Answer: At least 2 (one per class)}
\end{enumerate}



\subsection*{Pages 30--34 — Hard-Margin Constraints}

\textbf{Content extract}
\begin{itemize}
	\item Constraints:
	\[
	y_i(\mathbf{w}^\top \mathbf{x}_i + b) \ge 1
	\]
	\item Margin hyperplanes:
	\[
	\mathbf{w}^\top \mathbf{x} + b = \pm 1
	\]
\end{itemize}

\begin{enumerate}
	\item What does the constraint $y_i(\mathbf{w}^\top \mathbf{x}_i + b)\ge1$ ensure?
	\begin{enumerate}
		\item Maximum likelihood
		\item Correct classification with margin
		\item Zero training error only
		\item Kernel validity
	\end{enumerate}
	\textbf{Answer: (b)}
	
	\item True or False: All samples satisfy the constraint with equality.
	
	\textbf{Answer: False}
	
	\item Which samples satisfy equality?
	
	\textbf{Answer: Support vectors}
\end{enumerate}


\subsection*{Pages 35--42 — Primal Optimization}

\textbf{Content extract}
\begin{itemize}
	\item Maximizing margin is equivalent to minimizing $\frac{1}{2}\|\mathbf{w}\|^2$.
	\item Quadratic optimization with linear constraints.
\end{itemize}

\begin{enumerate}
	\item What objective function is minimized?
	\begin{enumerate}
		\item $\|\mathbf{w}\|$
		\item $\|\mathbf{w}\|^2$
		\item $\frac{1}{2}\|\mathbf{w}\|^2$
		\item $\|\mathbf{x}\|^2$
	\end{enumerate}
	\textbf{Answer: (c)}
	
	\item True or False: The primal problem is convex.
	
	\textbf{Answer: True}
	
	\item What type of optimization problem is this?
	
	\textbf{Answer: Quadratic programming}
\end{enumerate}


\subsection*{Pages 43--51 — Dual Problem and KKT Conditions}

\textbf{Content extract}
\begin{itemize}
	\item Lagrange multipliers $\alpha_i$ are introduced.
	\item Dual formulation depends only on dot products.
	\item KKT conditions identify support vectors.
\end{itemize}

\begin{enumerate}
	\item What do $\alpha_i > 0$ indicate?
	\begin{enumerate}
		\item Misclassified points
		\item Redundant points
		\item Support vectors
		\item Kernel values
	\end{enumerate}
	\textbf{Answer: (c)}
	
	\item True or False: All training points appear in the decision function.
	
	\textbf{Answer: False}
	
	\item How is $b$ computed?
	
	\textbf{Answer: Using any support vector (usually averaged)}
\end{enumerate}


\subsection*{Pages 52--58 — Linear SVM Example}

\textbf{Content extract}
\begin{itemize}
	\item Numerical example of linear SVM.
	\item Explicit computation of $\mathbf{w}$ and $b$.
	\item Classification using $f(\mathbf{x})$.
\end{itemize}

\begin{enumerate}
	\item What does $f(\mathbf{x})$ represent?
	\begin{enumerate}
		\item Probability
		\item Distance from hyperplane
		\item Classification score
		\item Kernel value
	\end{enumerate}
	\textbf{Answer: (c)}
	
	\item True or False: Larger $|f(\mathbf{x})|$ implies higher confidence.
	
	\textbf{Answer: True}
	
	\item If $f(\mathbf{x})<0$, what is the predicted class?
	
	\textbf{Answer: $-1$}
\end{enumerate}


\subsection*{Pages 59--70 — Soft-Margin SVM}

\textbf{Content extract}
\begin{itemize}
	\item Slack variables $\xi_i$ allow classification errors.
	\item Parameter $C$ controls margin–error tradeoff.
\end{itemize}

\begin{enumerate}
	\item What does a large $C$ emphasize?
	\begin{enumerate}
		\item Wider margin
		\item Fewer errors
		\item More regularization
		\item Kernel smoothness
	\end{enumerate}
	\textbf{Answer: (b)}
	
	\item True or False: $\xi_i>0$ implies a classification error or margin violation.
	
	\textbf{Answer: True}
	
	\item What happens when $C \to 0$?
	
	\textbf{Answer: Margin dominates, more errors allowed}
\end{enumerate}

\subsection*{Pages 71--78 — Overfitting and Generalization}

\textbf{Content extract}
\begin{itemize}
	\item Narrow margins lead to overfitting.
	\item Allowing errors improves generalization.
\end{itemize}

\begin{enumerate}
	\item What causes overfitting in SVM?
	\begin{enumerate}
		\item Large margin
		\item Small margin
		\item Kernel choice
		\item Feature scaling
	\end{enumerate}
	\textbf{Answer: (b)}
	
	\item True or False: Soft-margin SVM reduces overfitting.
	
	\textbf{Answer: True}
	
	\item Which parameter mainly controls overfitting?
	
	\textbf{Answer: $C$}
\end{enumerate}

\subsection*{Pages 79--95 — Nonlinear SVM and Kernels}

\textbf{Content extract}
\begin{itemize}
	\item Data may not be separable by a hyperplane.
	\item Kernel trick maps data implicitly to higher dimensions.
	\item Common kernels: linear, polynomial, RBF, sigmoid.
\end{itemize}

\begin{enumerate}
	\item What does the kernel trick avoid?
	\begin{enumerate}
		\item Optimization
		\item Explicit feature mapping
		\item Regularization
		\item Slack variables
	\end{enumerate}
	\textbf{Answer: (b)}
	
	\item True or False: Kernels compute dot products in feature space.
	
	\textbf{Answer: True}
	
	\item Which kernel corresponds to infinite-dimensional mapping?
	
	\textbf{Answer: Gaussian (RBF)}
\end{enumerate}


\subsection*{Pages 96--END — Multiclass Example}

\textbf{Content extract}
\begin{itemize}
	\item Iris dataset classification.
	\item One-vs-all strategy for multiclass SVM.
	\item Linear and polynomial kernels used.
\end{itemize}

\begin{enumerate}
	\item How is multiclass SVM implemented here?
	\begin{enumerate}
		\item One-vs-one
		\item One-vs-all
		\item Softmax
		\item Decision trees
	\end{enumerate}
	\textbf{Answer: (b)}
	
	\item True or False: Linear SVM separates setosa perfectly.
	
	\textbf{Answer: True}
	
	\item How many binary classifiers are needed for one-vs-all with 3 classes?
	
	\textbf{Answer: 3}
\end{enumerate}

\end{document}
